\section{Methods}
\label{sec:methods}

We report a single estimator under two sampling protocols. This section proceeds in the order of the pipeline. We first describe the dataset and the rotation-augmentation protocol that supplies the angle labels (\S\ref{sec:methods:data}). We then define the two sampling protocols and state what each one measures (\S\ref{sec:methods:protocols}). Next come the frozen-backbone feature extractor and the orientation-preserving grid-pooling head suited to an angular target (\S\ref{sec:methods:backbone}), the regression head and its Optuna tuning on cached features (\S\ref{sec:methods:head}), and the cross-validated ensembles that give our best estimator (\S\ref{sec:methods:ensemble}). Finally we describe a classical structure-tensor prior and its circular fusion with the learned estimate (\S\ref{sec:methods:fusion}) and the evaluation suite (\S\ref{sec:methods:eval}). Figure~\ref{fig:pipeline} summarizes the full data-to-evaluation pipeline.

% End-to-end experiment pipeline: data -> augmentation -> two sampling protocols
% -> cached frozen features -> per-protocol Optuna tuning -> ensemble -> evaluation.
% Self-contained TikZ (needs tikz + arrows.meta, positioning, fit libraries).
\begin{figure}[htbp]
\centering
\resizebox{\linewidth}{!}{%
\begin{tikzpicture}[
  font=\small,
  >={Stealth[length=2.2mm]},
  every node/.style={align=center},
  stage/.style={draw, rounded corners=2pt, minimum height=11mm, minimum width=20mm, fill=blue!5},
  data/.style={stage, fill=black!7},
  proto/.style={draw, rounded corners=2pt, minimum height=7mm, minimum width=24mm, fill=orange!14, font=\footnotesize},
  eval/.style={stage, fill=green!8, minimum width=24mm},
  node distance=8mm and 10mm,
]
  \node[data]  (raw)  {84 base\\images};
  \node[stage, right=of raw]  (aug)  {augment\\$\pm60^\circ/5^\circ$ + CLAHE};
  \node[stage, right=of aug]  (corp) {2{,}100-image\\corpus};
  \node[stage, right=of corp] (feat) {cached frozen\\features\\(1 / backbone)};
  \node[stage, right=of feat] (tune) {Optuna TPE\\tuning\\(per protocol)};
  \node[stage, right=of tune] (ens)  {5-backbone\\ensemble\\mean / stacked};
  \node[eval,  right=of ens]  (ev)   {evaluation};

  \draw[->] (raw)--(aug); \draw[->] (aug)--(corp); \draw[->] (corp)--(feat);
  \draw[->] (feat)--(tune); \draw[->] (tune)--(ens); \draw[->] (ens)--(ev);

  % the two sampling protocols feed the CV that scores tuning
  \node[proto, above=8mm of tune] (pimg) {image-level (random)};
  \node[proto, below=8mm of tune] (ppat) {patient-level (GroupKFold)};
  \draw[->, dashed, gray] (pimg) -- (tune);
  \draw[->, dashed, gray] (ppat) -- (tune);

  % what evaluation comprises
  \node[below=7mm of ev, text width=30mm, font=\footnotesize] (evlist)
       {conformal intervals,\\Bland--Altman, TTA,\\Grad-CAM,\\structure-tensor fusion};
  \draw[->, dashed, gray] (ev) -- (evlist);

  \begin{scope}[on background layer]
    \node[draw=black!25, dashed, rounded corners, fit=(pimg)(ppat),
          inner sep=2mm, label=left:{\footnotesize two protocols}] {};
  \end{scope}
\end{tikzpicture}}
\caption{Experiment pipeline. The 84 base images are expanded by a deterministic
rotation sweep (with CLAHE) into a $\sim$2{,}100-image corpus. Each backbone's
frozen features are extracted once and cached, then scored under \emph{two
complementary sampling protocols} - image-level (random) and patient-level
(grouped by volunteer) - with per-protocol Optuna TPE tuning of the head. The
five tuned backbones are combined (mean or stacked ensemble) and assessed with a
clinical-grade evaluation suite.}
\label{fig:pipeline}
\end{figure}


\subsection{Data and rotation augmentation}
\label{sec:methods:data}

The cohort is the one introduced in the original EMBC~2019 study \citep{Patil2019}, drawn from the public SPLab (Brno) ultrasound database. It comprises \num{84} B-mode images of the common carotid artery in longitudinal section, from ten volunteers (mean age $27.5\pm3.5$ years, usually eight images per volunteer). The subjects are young and healthy, the site is the right common carotid in longitudinal view, and acquisition used a single scanner; we draw no conclusions about diseased, aged, or tortuous vessels from this cohort. Images were acquired on a Sonix~OP scanner with two linear-array transducers (\num{10} and \SI{14}{\mega\hertz}, chosen for superficial vascular scanning). The acquisition protocol was fixed: subject supine, neck rotated to the left, right common carotid imaged at a resolution of roughly $390\times330$ pixels. The supervised target is the Doppler angle $\theta$, the angle between the ultrasound beam and the long axis of the vessel (Fig.~\ref{fig:angle_convention}), which enters the Doppler equation~\eqref{eq:doppler} through $\cos\theta$ and is set by hand at the scanner. Because the recovered velocity scales as $1/\cos\theta$, an error in $\theta$ is a leading, operator-dependent source of velocity error (\S\ref{sec:intro}). The database carries no clinical Doppler-angle annotation, so a reference $\theta$ was measured offline by one operator in a custom MATLAB interface, drawing a line along the vessel wall and recording its inclination to the image's vertical axis, with $\theta\in[0,\SI{180}{\degree}]$. This single hand-drawn reading is the only ground truth available per image; its own measurement error is unknown. The estimator reads $\theta$ directly from a single grayscale image, using neither color Doppler nor vessel segmentation.

% Doppler-angle convention: theta is the angle between the image's vertical axis
% (the ultrasound beam direction) and the vessel long axis, theta in [0,180].
% Self-contained TikZ (needs tikz + arrows.meta, already loaded in the preamble).
\begin{figure}[htbp]
\centering
\begin{tikzpicture}[
  font=\small,
  >={Stealth[length=2.4mm]},
  axis/.style={->, blue!55!black, line width=1.0pt},
  vessel/.style={line width=8pt, gray!30, line cap=round},
]
  % image frame
  \draw[gray!40, rounded corners=1pt] (-2.8,-2.5) rectangle (3.0,2.5);
  \node[gray!60, font=\footnotesize] at (1.75,-2.25) {B-mode frame};

  % the vessel as an undirected line at angle theta from the vertical
  \def\thetadeg{58}
  \draw[vessel] ({-2.4*sin(\thetadeg)},{-2.4*cos(\thetadeg)})
             -- ({2.4*sin(\thetadeg)},{2.4*cos(\thetadeg)});
  \draw[black, line width=0.9pt]
       ({-2.4*sin(\thetadeg)},{-2.4*cos(\thetadeg)})
    -- ({2.4*sin(\thetadeg)},{2.4*cos(\thetadeg)});
  \node[font=\footnotesize, anchor=south west]
       at ({2.0*sin(\thetadeg)},{2.0*cos(\thetadeg)}) {vessel wall};
  % faint flow arrow along the vessel (direction-agnostic; theta is undirected)
  \draw[->, red!55!black, line width=0.9pt]
       (0,0) -- ({1.15*sin(\thetadeg)},{1.15*cos(\thetadeg)});
  \node[red!55!black, font=\footnotesize]
       at ({1.15*sin(\thetadeg)+0.35},{1.15*cos(\thetadeg)-0.12}) {flow};

  % beam axis = image vertical, through the same intersection
  \draw[axis] (0,2.45) -- (0,-2.45);
  \node[blue!55!black, font=\footnotesize, anchor=south east] at (-0.12,1.9)
       {beam axis};
  \node[blue!55!black, font=\footnotesize, anchor=north east] at (-0.12,1.88)
       {(image vertical)};

  % theta: the angle between the upward vertical and the vessel (upper-right ray)
  \draw[->, black, line width=0.9pt]
       (0,1.0) arc[start angle=90, end angle={90-\thetadeg}, radius=1.0];
  \node at ({0.62*sin(\thetadeg/2)},{0.62*cos(\thetadeg/2)+0.32}) {$\theta$};
\end{tikzpicture}
\caption{Doppler-angle convention. The Doppler angle $\theta$ is the angle
between the ultrasound beam (the image's vertical axis) and the vessel long axis,
read off a single B-mode frame by aligning a line with the vessel wall and
measuring its inclination to the vertical; $\theta$ ranges over
$[0,\SI{180}{\degree}]$ and is direction-agnostic. The recovered velocity scales
as $1/\cos\theta$, so the accuracy of this angle directly controls the accuracy of
the reported velocity.}
\label{fig:angle_convention}
\end{figure}


To build a regression corpus with known labels, each base image is rotated through a fixed grid of angles spanning $[-\num{60},+\num{60}]\,\si{\degree}$ in \SI{5}{\degree} steps, yielding \num{25} oriented views per image and $\num{84}\times\num{25}=\num{2100}$ augmented views in total. The label of a rotated view is the base-image reading plus the applied rotation. The applied increment is exact, so the network learns to read off relative image orientation; the label is therefore exact only relative to the single base-image reading, and inherits whatever error that reading carries. Prior to feature extraction each view is contrast-equalized with Contrast-Limited Adaptive Histogram Equalization (CLAHE) \citep{Zuiderveld1994}, which normalizes the wide dynamic range of B-mode speckle and stabilizes the frozen-backbone features across acquisitions. This rotation-augmentation recipe is a standard way to build a labelled regression corpus from a small set of base images. The \num{2100} views are not \num{2100} independent observations: they are \num{25} geometrically coupled copies of each of \num{84} base images, drawn from ten volunteers, giving a three-level hierarchy (volunteer $\rightarrow$ base image $\rightarrow$ rotated view) that the two sampling protocols below treat differently. Figure~\ref{fig:augmentation} shows the rotation sweep for one base image, each view labelled with its resulting angle.

\begin{figure}[htbp]
  \centering
  \includegraphics[width=0.86\linewidth]{figure2_augmentation.png}
  \caption{Rotation-augmentation sweep for a single base image: the image is rotated through $[-60,+60]\si{\degree}$ in \SI{5}{\degree} steps (\num{25} views), and the applied rotation defines each view's exact Doppler-angle label (panel titles). CLAHE contrast equalization is applied before feature extraction. This manufactures a labelled regression corpus of $\num{84}\times\num{25}=\num{2100}$ views from the \num{84} base images.}
  \label{fig:augmentation}
\end{figure}

\subsection{Two sampling protocols}
\label{sec:methods:protocols}

The hierarchical structure of the corpus admits more than one principled train/test partition. We report and tune the estimator under two complementary sampling protocols, each answering a distinct question:

\begin{description}[leftmargin=1.25em,style=nextline]
  \item[Image-level sampling :] The \num{2100} rotated views are partitioned at random over the augmented corpus. This is the original study's protocol. It measures angle accuracy across the full population of orientations and imaging conditions the augmentation spans, the regime in which the published headline numbers were reported. Views of one base image can fall on both sides of the split, so this protocol does not test transfer to a new subject.
  \item[Patient-level sampling :] Whole subjects are held out with a grouped split (\textsc{GroupKFold}), so neither a test subject's rotated views nor any of their other images appear in training. The SPLab filenames do not encode volunteer identity, so we recover subject groups by clustering on acquisition time. This yields \num{12} patient-proxy groups over the \num{84} images, more than the ten volunteers, so the proxy grouping is a heuristic and we have no ground-truth subject labels to validate it against; the patient-disjoint guarantee is only as good as that clustering. These \num{12} groups are the unit of all patient-level cross-validation and per-patient aggregation in this paper. The groups are markedly imbalanced (from a single base image up to \num{21}, i.e.\ \num{25}--\num{525} rotated views), so one dominant group can land in a single fold. This inflates per-fold variance and is reflected in the large patient-level standard deviations. The protocol is a stricter lens: it measures generalization to a previously unseen patient. That is harder, because the estimator can no longer rely on subject-specific appearance.
\end{description}

The two protocols answer different questions, and we report both. Holding the model, backbone, head, and training budget fixed while changing only the sampling protocol gives an informative measurement: the patient-level figure quantifies how much of the within-population accuracy is anatomy-specific. Across backbones the patient-level MAPE is roughly double the image-level MAPE, which is what one expects from a small, single-center cohort. The patient-level numbers are the clinically relevant ones, and we read them as such. We tune each protocol separately and report both throughout.

\subsection{Frozen backbones and the grid-pooling head}
\label{sec:methods:backbone}

We use ImageNet-pretrained \citep{Deng2009} convolutional backbones as \emph{frozen} feature extractors, with the classification top removed (\texttt{include\_top=False}); no backbone weights are updated. We evaluate five architectures from the classical ImageNet generation: VGG19 \citep{Simonyan2015}, ResNet50 \citep{He2016}, DenseNet201 \citep{Huang2017}, Xception \citep{Chollet2017}, and InceptionV3 \citep{Szegedy2016}. Freezing keeps the comparison faithful to the compute regime of the original work. It also makes the feature maps cacheable, so the head search of \S\ref{sec:methods:head} runs on precomputed vectors. Figure~\ref{fig:architecture} shows the architecture.

% Estimator architecture: frozen backbone -> orientation-preserving grid pooling
% -> shallow regression head -> angle. Self-contained TikZ (needs tikz +
% arrows.meta, positioning libraries, loaded in the preamble).
\begin{figure}[htbp]
\centering
\resizebox{\linewidth}{!}{%
\begin{tikzpicture}[
  font=\small,
  >={Stealth[length=2.2mm]},
  every node/.style={align=center},
  box/.style={draw, rounded corners=2pt, minimum height=11mm, minimum width=17mm, fill=blue!4},
  frozen/.style={box, fill=black!8},
  pool/.style={box, fill=orange!16},
  head/.style={box, fill=green!8},
  io/.style={box, fill=blue!10, minimum width=13mm},
  node distance=7mm,
]
  \node[io]      (img)  {B-mode\\image\\$H{\times}W$};
  \node[frozen, right=of img]  (bb)  {frozen ImageNet\\backbone\\\texttt{include\_top=False}};
  \node[box,    right=of bb]   (fmap){conv map\\$h{\times}w{\times}c$};
  \node[pool,   right=of fmap] (grid){\textbf{grid pool}\\$G{\times}G$};
  \node[box,    right=of grid] (flat){flatten\\$G^2c$};
  \node[head,   right=of flat] (head){BN\,$\to$\,Dense\\$\to$\,ReLU\,$\to$\,Drop\\$\to$\,Dense};
  \node[io,     right=of head] (out) {$\hat\theta$\\(degrees)};

  \draw[->] (img)--(bb); \draw[->] (bb)--(fmap); \draw[->] (fmap)--(grid);
  \draw[->] (grid)--(flat); \draw[->] (flat)--(head); \draw[->] (head)--(out);

  % the load-bearing contrast, annotated under the pooling stage
  \node[below=5mm of grid, text width=42mm, font=\footnotesize] (note)
       {keeps coarse spatial layout\\($G{\times}G$, orientation-aware)\\[1pt]
        \emph{vs} global avg-pool: $1{\times}1$,\\approximately rotation-invariant};
  \draw[->, dashed, gray] (note) -- (grid);

  \begin{scope}[on background layer]
    \node[draw=black!25, dashed, rounded corners, fit=(bb)(fmap), inner sep=2.5mm,
          label=above:{frozen feature extractor}] {};
    \node[draw=black!25, dashed, rounded corners, fit=(grid)(flat)(head), inner sep=2.5mm,
          label=above:{trained head}] {};
  \end{scope}
\end{tikzpicture}}
\caption{Estimator architecture. A frozen ImageNet backbone produces a
convolutional feature map; \textbf{orientation-preserving grid pooling} averages
that map to a small $G{\times}G$ grid (rather than a single $1{\times}1$ vector)
so the coarse spatial layout that encodes vessel orientation survives. The
flattened grid feeds a shallow regression head that outputs the Doppler angle
$\hat\theta$. Only the head is trained.}
\label{fig:architecture}
\end{figure}


The main design choice is how the final spatial feature map is collapsed into a fixed-length vector. The conventional global average pooling (GAP) head averages activations over all spatial locations and is, by construction, approximately rotation-invariant, so it discards the orientation signal the target depends on. That is the wrong inductive bias for an angular target. We replace GAP with an orientation-preserving grid-pooling head: the final feature map is average-pooled onto a small $g\times g$ spatial grid and the per-cell descriptors are flattened and concatenated, so the spatial location of each activation is retained. The frozen backbone already encodes the vessel orientation, and the work is in not pooling it away.

Restoring this coarse spatial structure roughly halves the error. Under image-level five-fold cross-validation, frozen DenseNet201 improves from \SI{10.85}{\percent} MAPE with a GAP head to \SI{4.58}{\percent} with grid pooling. On the original single augmented split it reaches \SI{5.84}{\percent} MAPE (\SI{3.77}{\degree} MAE, $R^2=0.982$), which is the protocol and the number reported in Table~I of the EMBC~2019 paper; this single-split figure is the one that replicates the published result, while the \SI{4.58}{\percent} cross-validated figure is a different estimate under a different protocol. Table~\ref{tab:replication} reports the replication across all five frozen backbones, with DenseNet201 the strongest.

A frozen bake-off under patient-level five-fold cross-validation supports the backbone choice, with the caveat that it rests on \num{84} base images and the same large fold variance noted above, so the ordering should be read against the reported standard deviations rather than as a clean ranking. On the extraction-matched comparison, frozen DenseNet201 (\SI{14.1}{\percent} MAPE) edges out the best modern backbone, ConvNeXt-Base \citep{Liu2022} (\SI{15.7}{\percent} MAPE), with the EfficientNet \citep{Tan2019} family at $\sim\!17$--\SI{21}{\percent} MAPE (Sec.~\ref{subsec:bakeoff}). On this narrow domain, no newer or larger backbone clearly improves on the 2017 DenseNet. DenseNet201, which gives the best $R^2$, wins the replication, and remains strongest after tuning, is carried forward as the primary backbone.


\subsection{Regression head and Optuna tuning on cached features}
\label{sec:methods:head}

On top of the grid-pooled features we place a compact head, \texttt{BatchNorm} $\rightarrow$ \texttt{Dense} $\rightarrow$ \texttt{ReLU} $\rightarrow$ \texttt{Dropout} $\rightarrow$ \texttt{Dense}, that outputs the scalar angle estimate. All models are trained with Adam \citep{Kingma2015} at base learning rate \num{1e-4} under a mean-squared-error loss; the backbone stays frozen throughout. Because the frozen, grid-pooled features can be computed once and cached, head training and the entire hyperparameter search operate on these cached vectors, which makes a large search affordable without a GPU.

We tune the head with the Tree-structured Parzen Estimator (TPE) sampler \citep{Bergstra2011} as implemented in Optuna \citep{Akiba2019}, searching over head depth and width, dropout rate, $L_2$ weight decay, whether batch normalization is applied, learning rate, batch size, and early-stopping patience. Each protocol is tuned separately: every objective evaluation uses the same sampling protocol it reports, scored on that protocol's five-fold cross-validation, so the image-level and patient-level estimators are each optimized for the question they answer. A single feature extraction per backbone serves both protocols. Because the search selects hyperparameters on the same five folds that are then reported, the tuned single-model figures are optimistic relative to a held-out set: they are in-sample to the search, as no nested or outer test fold is used. Tuning moves the single-model DenseNet201 result, with frozen and tuned both scored by five-fold cross-validation (mean over folds), from \SI{4.58}{\percent} to \SI{4.03}{\percent} MAPE (\SI{3.00}{\degree} MAE) under image-level sampling, and from \SI{12.59}{\percent} to \SI{10.80}{\percent} MAPE (\SI{8.62}{\degree} MAE, $R^2=0.886$) under patient-level sampling. The per-model gains are small but consistent, and they calibrate the five members enough that a plain mean ensemble becomes effective.

\subsection{Cross-validated ensembles}
\label{sec:methods:ensemble}

We combine the tuned heads over all five backbones. Two combiners are reported: a plain mean of the member predictions, and a stacked generalization \citep{Wolpert1992} in which a Ridge meta-learner is fit on the member outputs. The stacker and every reported ensemble figure are computed from out-of-fold (OOF) predictions, pooled across folds into a single metric: under five-fold cross-validation each member predicts only on the fold it did not train on, so under patient-level sampling every ensemble number is evaluated on patients absent from the corresponding fold's training. This pooled-OOF aggregation differs from the mean-over-folds aggregation used for the single-model figures of \S\ref{sec:methods:head}, so single-model and ensemble metrics are not strictly the same estimator and should be compared with that in mind. The Ridge meta-learner is fit on OOF member outputs over only \num{12} patient groups, which leaves the meta-layer at some risk of overfitting; we did not run a separate nested OOF loop for the stacker. The mean ensemble reaches \SI{3.03}{\percent} MAPE (\SI{2.09}{\degree} MAE, $R^2=0.994$) under image-level sampling and \SI{9.89}{\percent} MAPE (\SI{6.89}{\degree} MAE, $R^2=0.932$) under patient-level sampling. The stacked ensemble is the estimator we carry forward, reaching \SI{2.79}{\percent} MAPE (\SI{1.96}{\degree} MAE, $R^2=0.995$) under image-level sampling and \SI{8.53}{\percent} MAPE (\SI{5.93}{\degree} MAE, $R^2=0.952$) under patient-level sampling. Ensembling gives the largest single gain we observe: combining the five backbones recovers accuracy that no single model reaches, consistent with the members making partly independent errors, though we do not report an error-correlation analysis to quantify that independence. The patient-level OOF predictions are also the substrate for the calibrated reporting of \S\ref{sec:methods:eval}.

\subsection{Classical structure-tensor prior and circular fusion}
\label{sec:methods:fusion}

To situate the learned estimator against a transparent geometric baseline, we also compute a purely classical, image-only orientation prior from the structure tensor: the dominant local-gradient orientation, aggregated over the vessel region, gives an angle estimate with no learning and no training data. Evaluated on the narrow band of base-image angles, this prior reaches \SI{3.16}{\degree} MAE. Because an orientation is defined only up to a half-turn, we fuse the learned and classical estimates in the circular $2\theta$ domain, averaging $(\cos 2\theta,\sin 2\theta)$ before mapping back. On the same base-image set the fused estimate reaches \SI{2.72}{\degree} MAE, below either component alone. This comparison is on \num{84} base images and a different, smaller evaluation set than the rotation-augmented corpus used elsewhere; we read the fused gain as suggestive of partly complementary information between the learned model and the geometric cue, not as an established decomposition, since we report no residual correlations or confidence intervals.

\subsection{Evaluation suite}
\label{sec:methods:eval}

Beyond a point estimate, we report calibrated diagnostics, all derived from the patient-level OOF predictions of the tuned DenseNet201. We construct distribution-free, finite-sample prediction intervals by split conformal regression \citep{Vovk2005,Lei2018}, using a patient-disjoint calibration/test partition so that calibration residuals come from subjects absent from both the model's training and the reported test set. Marginal conformal coverage is a guarantee in expectation; on a cohort of \num{12} imbalanced groups the calibration set is small, so the coverage actually observed on the single test partition is itself noisy, and the achieved coverage and calibration-set size should be read alongside the interval width in \S\ref{sec:evaluation}. Agreement with the reference angle is summarized by a Bland--Altman analysis \citep{BlandAltman1986}, reporting bias and limits of agreement. Because only a single reference reading (the MATLAB-GUI angle) is available per image, this is a method-versus-reference comparison, not an inter-observer study, and the precision of the reference is unknown. At inference we apply rotation test-time augmentation, predicting each base image over its rotated views and aggregating by the circular median, which reduces base-image MAE; note that this reuses the same rotation family that generated the labels, so the gain is measured over transformations of training-adjacent data rather than independent views. Finally, Grad-CAM \citep{Selvaraju2017} attributions are computed on the frozen backbone as a qualitative check that the estimator keys on the vessel wall, the anatomy that defines the flow axis, rather than on background structure; this is a saliency inspection, not a formal verification. The full clinical evaluation is reported in \S\ref{sec:evaluation}.
